\subsection{线性调制系统的抗噪声性能}
\subsubsection{含噪解调模型和性能指标}
带限 AWGN 信道和解调模型如下
\begin{equation*}
    s_m(t)+n(t)\rightarrow BPF\rightarrow s_i(t), n_i(t)\rightarrow \text{解调器}\rightarrow m_o(t), n_o(t)
\end{equation*}
其中, $n(t)$ 是高斯白噪声, 经过 BPF (降低噪声带宽, 降低噪声频率) 后转为 $n_i(t)$ 窄带高斯噪声. 关于窄带随机过程的特性见 \ref{3.5}.

\textbf{输出信噪比}
\begin{equation}
    \frac{S_o}{N_o}=\frac{\overline{m_o(t)^2}}{n_o(t)^2}.
\end{equation}

\textbf{信噪比增益}
\begin{equation}
    G=\frac{S_o/N_o}{S_i/N_i}.
\end{equation}

\subsubsection{DSB 相干解调系统的抗噪声性能}
容易得到
\begin{gather}
    s_i(t)=m(t)\cos\omega_ct,\quad m_o(t)=\frac{1}{2}m(t); \\
    S_i=\overline{s_i^2(t)}=\frac{1}{2}\overline{m^2(t)},\quad S_o=\overline{m_o^2(t)}=\frac{1}{4}\overline{m^2(t)}; \\
    n_o(t)=\frac{1}{2}n_i(t).
\end{gather}

因此
\begin{gather}
    \frac{S_o}{N_o}=\frac{\overline{m^2(t)}}{n_0B}, \\
    G_{DSB}=2.
\end{gather}

相干解调将噪声中的正交分量滤除, 使信噪比改善.

\subsubsection{SSB 相干解调系统的抗噪声性能}
同理
\begin{gather}
    \frac{S_o}{N_o}=\frac{\overline{m^2(t)}}{4n_0B}, \\
    G_{SSB}=1.
\end{gather}

相干解调将噪声和调制信号的正交分量都有滤除.

在相同的 $S_i$, $n_0$, $f_H$ ($B_{DSB}=2f_H$, $B_{SSB}=f_H$) 下, DSB 和 SSB 抗噪声性能相同.

\subsubsection{AM 包络检波系统的抗噪声性能}
信噪比很大时,
\begin{gather}
    S_o=\overline{m^2(t)}, \quad N_o=N_i=\overline{n(t)}; \\
    G_{AM}=\frac{2\overline{m^2(t)}}{A_0^2+\overline{m^2(t)}}\leq\frac{2}{3}. \label{eq:5.2 G AM}
\end{gather}

信噪比很小时, 由于\textbf{门限效应}, 信号被扰乱成噪声, 输出信噪比急剧恶化.

\begin{exampleprob}
    发射信号为 $s_m(t)=[A_0+m(t)]\cos(\omega_ct)$, 其中 $f_H=5\ \text{kHz}$, $\overline{m^2(t)}=500\ \text{W}$, $f_c=100\ \text{kHz}$, 载波功率为 1000 W, 噪声单边功率谱密度为 $10^{-10}\ \text{W/Hz}$, 信道衰减 70 dB, BPF 高为 1 宽为 $B$.
    \begin{enumerate}[(1)]
        \item 确定 BPF 中心频率 $f_0$ 和带宽 $B$;
        \item 计算解调器输入信噪比;
        \item 计算解调器输出信噪比和调制制度增益.
    \end{enumerate}

    \begin{solution}
        显然这是 AM 调制.

        \begin{enumerate}[(1)]
            \item $f_0=f_c=100\ \text{kHz}$, $B=2f_H=10\ \text{kHz}$.
            \item 实际发射信号功率为
                  \begin{equation*}
                      S_T=\frac{A_0^2}{2}+\frac{\overline{m^2(t)}}{2}=1000+250=1250\ \text{W}.
                  \end{equation*}

                  因为信道存在衰减, 实际接收信号功率为
                  \begin{equation*}
                      S_i=\frac{S_T}{70\ \text{dB}}=125\ \mu\text{W}.
                  \end{equation*}

                  噪声功率为
                  \begin{equation*}
                      N_i=n_0B=1\ \mu\text{W}.
                  \end{equation*}

                  所以
                  \begin{equation*}
                      SNR_i=125.
                  \end{equation*}
            \item 因为输入信噪比很大, 所以可以使用式 (\ref{eq:5.2 G AM}): $G_{AM}=0.4$. 所以输出信噪比为
                  \begin{equation*}
                      SNR_o=G_{AM}SNR_i=50.
                  \end{equation*}
        \end{enumerate}
    \end{solution}
\end{exampleprob}
